سیستم به صورت زیر تعریف شده است:
$$
T\{x(t)\}=y(t)=x(2)+x(t-6).
$$
دو ورودی دلخواه $x_1(t),x_2(t)$ و ضرایب $a,b$ را در نظر بگیرید:
\begin{flalign*}
T\{ax_1+bx_2\}
&=(ax_1+bx_2)(2) + (ax_1+bx_2)(t-6) \\
&=a x_1(2)+b x_2(2) + a x_1(t-6)+b x_2(t-6) \\
&=a\big(x_1(2)+x_1(t-6)\big)+b\big(x_2(2)+x_2(t-6)\big) \\
&=aT\{x_1\}+bT\{x_2\}. &&
\end{flalign*}
پس سیستم خطی است.

برای علی بودن باید $y(t_0)$ فقط به $x(\tau)$ برای $\tau \le t_0$ وابسته باشد.
اما
$$
y(t_0)=x(2)+x(t_0-6).
$$
ترم $x(2)$ همیشه به مقدار ورودی در زمان $\tau=2$ نیاز دارد. اگر $t_0<2$ باشد، آنگاه $\tau=2>t_0$ و خروجی به ورودیِ آینده وابسته می‌شود.
پس سیستم غیرعلی است.

ورودی شیفت‌داده‌شده را $x_1(t)=x(t-t_0)$ بگیرید. خروجی متناظر:
\begin{flalign*}
y_1(t)=T\{x_1\}
&=x_1(2)+x_1(t-6) \\
&=x(2-t_0)+x(t-6-t_0). &&
\end{flalign*}
از طرفی اگر ابتدا خروجی $y(t)$ را شیفت دهیم:
\begin{flalign*}
y(t-t_0)=x(2)+x(t-t_0-6). &&
\end{flalign*}
به طور کلی $x(2-t_0)\neq x(2)$، پس
$$
y_1(t)\neq y(t-t_0)
$$
و سیستم زمان‌متغیر است (زمان‌ناوردا نیست).

اگر ورودی کراندار باشد، یعنی $\lvert x(t)\rvert \le M$ برای همه‌ی $t$، آنگاه:
\begin{flalign*}
|y(t)|
&=|x(2)+x(t-6)|
\le |x(2)|+|x(t-6)|
\le M+M=2M. &&
\end{flalign*}
پس خروجی نیز کراندار است و سیستم \lr{BIBO stable} می‌باشد.